\chapter{Algorithms}
\label{chap:algorithms}

\chapterintro{An algorithm is a method for solving a problem. This chapter introduces algorithmic thinking using simple examples before later applying it to machine learning and LLM systems.}

\learningobjectives{
\begin{itemize}
    \item Define an algorithm.
    \item Explain step-by-step problem solving.
    \item Understand inputs, outputs, and procedures.
    \item Write simple algorithms in plain English and Python.
\end{itemize}}

\section{What Is an Algorithm?}

\begin{definition}
An algorithm is a finite sequence of well-defined steps for solving a problem or completing a task.
\end{definition}

A recipe is an everyday algorithm. It starts with ingredients, follows steps, and produces a meal.

\section{Inputs and Outputs}

An algorithm receives inputs and produces outputs.

For example, an algorithm for adding two numbers receives two numbers and returns their sum.

\begin{lstlisting}
def add(a: float, b: float) -> float:
    return a + b
\end{lstlisting}

\section{A Plain-English Algorithm}

Algorithm: find the largest number in a list.

\begin{enumerate}
    \item Start with the first number as the largest seen so far.
    \item Look at the next number.
    \item If it is larger, remember it as the new largest.
    \item Repeat until no numbers remain.
    \item Return the largest number.
\end{enumerate}

\section{Python Implementation}

\begin{lstlisting}
from collections.abc import Sequence


def maximum(values: Sequence[float]) -> float:
    if not values:
        raise ValueError("values cannot be empty")

    largest = values[0]
    for value in values[1:]:
        if value > largest:
            largest = value
    return largest


print(maximum([3, 1, 9, 2]))
\end{lstlisting}

\section{Correctness}

An algorithm is correct if it produces the required output for every valid input.

Correctness is important in AI engineering because hidden bugs can corrupt datasets, evaluations, and production decisions.

\section{Efficiency}

Efficiency describes how much time and memory an algorithm uses. Later we will study Big-O notation. For now, notice that some methods require more work than others.

\section{Algorithms in LLM Engineering}

LLM systems use algorithms for:

\begin{itemize}
    \item tokenization,
    \item embedding search,
    \item ranking,
    \item batching,
    \item caching,
    \item sampling,
    \item evaluation.
\end{itemize}

\section{Exercises}

\begin{exercise}
Write a plain-English algorithm for making a cup of tea.
\end{exercise}

\begin{solution}
A valid solution should describe finite, ordered steps such as boil water, place tea in cup, add water, wait, remove tea, and serve.
\end{solution}

\begin{exercise}
Write a Python function that returns the smallest number in a non-empty list.
\end{exercise}

\section{Portfolio Milestone}

Create an \texttt{algorithms-lab} folder containing Python implementations for maximum, minimum, average, and counting words in a sentence.

\section{Summary}

Algorithms are precise procedures for solving problems. They are central to all computing and appear throughout AI, from data cleaning to transformer inference.
